% Part: incompleteness
% Chapter: representability-in-q
% Section: composition-representable

\documentclass[../../../include/open-logic-section]{subfiles}

\begin{document}

\olfileid{inc}{req}{cmp}
\olsection{Komposisi Representabel ing $\Th{Q}$}

Upamana $h$ ditegesi dening
\[
h(x_0,\dots,x_{l-1}) = f(g_0(x_0,\dots,x_{l-1}), \dots,
g_{k-1}(x_0,\dots,x_{l-1})).
\]
ing ngendi kita wis nemokake !!{formula}s $!A_f, !A_{g_0}, \dots,
!A_{g_{k-1}}$ kang merepresentasikake fungsi $f$, lan $g_0$,
\dots,~$g_{k-1}$ miturut urutane. Kita kudu nemokake !!a{formula}~$!A_h$
kang merepresentasikake $h$.

Ayo wiwiti saka kasus prasaja, nalika kabeh fungsi nduweni $1$ argumen,
yaiku $h(x) = f(g(x))$. Yen $!A_f(y, z)$ merepresentasikake~$f$, lan
$!A_g(x, y)$ merepresentasikake~$g$, kita mbutuhake !!a{formula}~$!A_h(x, z)$
kang merepresentasikake~$h$. Gatekna yen $h(x) = z$ yen lan mung yen ana~$y$
saengga $z = f(y)$ lan $y = g(x)$. (Yen $h(x) = z$, mula $g(x)$ iku
$y$ kaya mangkono; yen $y$ kaya mangkono ana, mula amarga $y = g(x)$
lan $z = f(y)$, kita nduweni $z = f(g(x))$.) Iki nuduhake yen
$\lexists[y][(!A_g(x, y) \land !A_f(y,
  z))]$ calon kang cocog kanggo~$!A_h(x, z)$. Kita mung kudu mriksa yen
$\Th{Q}$ mbuktekake !!{formula}s kang dibutuhake.

\begin{prop}
\ollabel{prop:rep1}
Yen $h(n) = m$, mula $\Th{Q} \Proves !A_h(\num{n}, \num{m})$.
\end{prop}

\begin{proof}
Upamana $h(n) = m$, yaiku $f(g(n)) = m$. Tetepna $k = g(n)$. Mula
\begin{align*}
  \Th{Q} & \Proves !A_g(\num{n}, \num{k})
  \intertext{amarga $!A_g$ merepresentasikake~$g$, lan}
  \Th{Q} & \Proves !A_f(\num{k}, \num{m})
  \intertext{amarga $!A_f$ merepresentasikake~$f$. Dadi,}
  \Th{Q} & \Proves !A_g(\num{n}, \num{k}) \land !A_f(\num{k}, \num{m})
  \intertext{lan mula uga}
  \Th{Q} & \Proves \lexists[y][(!A_g(\num{n}, y) \land !A_f(y, \num{m}))],
\end{align*}
yaiku $\Th{Q} \Proves !A_h(\num{n}, \num{m})$.
\end{proof}

\begin{prop}
\ollabel{prop:rep2}
Yen $h(n) = m$, mula $\Th{Q} \Proves \lforall[z][(!A_h(\num{n}, z) \lif
  z = \num{m})]$.
\end{prop}

\begin{proof}
Upamana $h(n) = m$, yaiku $f(g(n)) = m$. Tetepna $k = g(n)$. Mula
\begin{align*}
  \Th{Q} & \Proves \lforall[y][(!A_g(\num{n}, y) \lif \eq[y][\num{k}])]
  \intertext{amarga $!A_g$ merepresentasikake~$g$, lan}
  \Th{Q} & \Proves \lforall[z][(!A_f(\num{k}, z) \lif \eq[z][\num{m}])]
  \intertext{amarga $!A_f$ merepresentasikake~$f$. Kanthi logika
    sithik wae, kita uga bisa nuduhake}
  \Th{Q} & \Proves \lforall[z][(\lexists[y][(!A_g(\num{n}, y) \land
      !A_f(y, z))] \lif \eq[z][\num{m}])].
\end{align*}
yaiku $\Th{Q} \Proves \lforall[y][(!A_h(\num n, y) \lif \eq[y][\num m])]$.
\end{proof}

Gagasan kang padha lumaku kanggo kasus luwih ruwet nalika $f$ lan~$g_i$
nduweni luwih saka~$1$ argumen.

\begin{prop}
\ollabel{prop:rep-composition}
Yen $!A_f(y_0, \dots, y_{k-1}, z)$ merepresentasikake $f(y_0, \dots, y_{k-1})$
ing~$\Th{Q}$, lan $!A_{g_i}(x_0, \dots, x_{l-1}, y)$ merepresentasikake
$g_i(x_0, \dots, x_{l-1})$ ing~$\Th{Q}$, mula
\begin{multline*}
  \lexists[y_0\dots][\lexists[y_{k-1}][(!A_{g_0}(x_0,\dots,x_{l-1},y_0) \land
      \dots \land {}]]\\
  !A_{g_{k-1}}(x_0,\dots,x_{l-1},y_{k-1}) \land !A_f(y_0,\dots,y_{k-1},z))
\end{multline*}
merepresentasikake
\[
h(x_0, \dots, x_{l-1}) = f(g_0(x_0, \dots, x_{l-1}), \dots, g_{k-1}(x_0,
\dots, x_{l-1})).
\]
\end{prop}

\begin{proof}
Latihan.
\end{proof}

\begin{prob}
Kanthi ngetutake bukti \olref[inc][req][cmp]{prop:rep2} lan
\olref[inc][req][cmp]{prop:rep2} minangka tuntunan, buktekna
\olref[inc][req][cmp]{prop:rep-composition} kanthi rinci.
\end{prob}

\end{document}
